% Source: sources/returns/S6_04_MAPS_MANIFOLDS/sections/manifolds.tex:129–296
\paragraph{Theorem D.1.4 (Proved: the integral homology of the universal cover).}
For every admissible triple, with $d=|p|$, one has
\[
 H_i(\widetilde X_p;\mathbb Z)=
 \begin{cases}
 \mathbb Z,&i=0,6,\\
 \mathbb Z^{d-1}\oplus\mathbb Z/d,&i=2,\\
 \mathbb Z/d,&i=3,\\
 \mathbb Z^{d-1},&i=4,\\
 0,&i=1,5.
 \end{cases}
\]
Thus for $d>1$ the cover is also not a connected sum of copies of
$S^2\times S^4$. In particular a torsion-free classification theorem
cannot be applied to these covers.
\begin{proof}
We give the integral attachment calculation. Write
$\gamma'=\gamma/d$ on the covering lattice
$\Lambda_d=\mathbb Z(d\widehat\gamma)\oplus\mathbb Z\widehat u
\oplus\mathbb Z\widehat w\oplus\mathbb Z\widehat\delta$.
Choose integers $h_j$ with
$m_jh_j\equiv\ell_j\pmod d$. Replace each meridian lift by its
product with the translation $-h_j\widehat\gamma$. Its $m_j$th power
has translation vector
$v_j-h_j\mathcal N_j\widehat\gamma$, where
\[
 \mathcal N_1\widehat\gamma=3\varepsilon,
 \qquad \mathcal N_2\widehat\gamma=4\varepsilon'.
\]
The non-invariant part of an affine translation is removed by
conjugation by a real torus translation: averaging projects to the
invariant part, and its complementary part is in the image of $A_j-I$
over $\mathbb R$. This produces precisely the same smooth local
quotient with invariant translation vector. Its new integers are
\[
 a_j=(\ell_j-m_jh_j)/d,\qquad
 a_0=(\ell_0-h_1-h_2)/d.
\]
The numerator defining $a_0$ is divisible by $d$ by the original
meridian relation. They satisfy
\[
 p'=12a_0-4a_1-3a_2=p/d\in\{1,-1\}.
\]
In the basis of $\Lambda_d$, the two matrices are
\[
 A_1^{(d)}=\begin{pmatrix}1&0&0&0\\6d&0&1&0\\-6d&-1&-1&0\\-2d&1&0&1\end{pmatrix},
 \quad
 A_2^{(d)}=\begin{pmatrix}1&0&0&0\\0&0&-1&0\\-6d&1&0&0\\3d&0&1&1\end{pmatrix}.
\]
Their coinvariant lattices are the same saturated rank-two quotients
as for $d=1$. The invariant forms are
$\psi_1=2u+w+3\delta$, $\psi_2=u+w+2\delta$,
and their common invariant alternating form is
$q_d=u\wedge w+6d\gamma'\wedge\delta$.

Let $\widetilde X^\circ$ be the complement of the cusp tube. The
finite reduced fibres, denoted $S'_1,S'_2$, have torsion-free first
homology: their abelianized affine presentation is the free rank-two
coinvariant lattice with one meridian adjoined, modulo
$m_ja_j^{\rm mer}=[v'_j]$, a primitive relation since
$\gcd(a_j,m_j)=1$. Poincare duality, Euler characteristic zero and
the universal coefficient theorem give $H_2(S'_j)=\mathbb Z^2$.
The intersection form is even and unimodular: the canonical class is
torsion for a bielliptic surface and $H^2(S'_j)$ is torsion free,
so its mod-two reduction, the second Stiefel--Whitney class, is zero.

The pullback images in $H^2(F_d)$ have respective integral bases
\[
 \begin{aligned}
 L_1&=\langle\gamma'\psi_1,\ q_d-2d\gamma'\psi_1\rangle,\\
 L_2&=\langle2\gamma'\psi_2,\ q_d-3d\gamma'\psi_2\rangle.
 \end{aligned}
\]
Here is a verification of the sublattice assertion, including its
index. The invariant lattices are
$\langle\gamma'\psi_j,q_d\rangle$, with intersection determinants
$-9$ and $-4$. Pullback multiplies intersection pairings by $3$ and
$4$, so its indices are respectively one and two. At the second
fibre, $q_d$ cannot descend since its square divided by four is
$3d$, odd. The class $\gamma'\psi_2$ cannot descend either: on the
torus swept by $\delta$ and the affine meridian its integral would
be $-a_2/2$, which is not integral. Of the three index-two sublattices
of the invariant lattice, only the parity sublattice
$\{a\gamma'\psi_2+bq_d:a\equiv b\pmod2\}$ excludes both classes.
The displayed basis generates that sublattice because $d$ is odd.
Its pairing divided by four, and the first pairing divided by three,
are the hyperbolic plane.

In homology bases dual to these cohomology bases, the map from the
common regular fibre to the two finite pieces is
\[
 \alpha_2=
 \begin{pmatrix}
 2&1&3&0&0&0\\
 -4d&-2d&0&1&0&0\\
 -2&-2&-4&0&0&0\\
 3d&3d&0&-1&0&0
 \end{pmatrix}.
\]
Its rank is three. Among its maximal minors are $2$, $3d$, $4d$;
their greatest common divisor is one because $d$ is odd. Therefore
its cokernel is $\mathbb Z$, whose primitive dual is
\[
 \Phi=(4d,2,3d,2).
\]
The degree-one map has kernel
$\mathbb Z(-\widehat u-\widehat w+\widehat\delta)$, exactly as at
$d=1$. Mayer--Vietoris for the two finite pieces therefore gives
\[
 0\longrightarrow\mathbb Z\longrightarrow H_2(\widetilde X^\circ)
 \longrightarrow\mathbb Z\longrightarrow0.
\]
Furthermore $H_1(\widetilde X^\circ)=\mathbb Z$, with primitive dual
$\Theta$ restricting to $12\gamma'$ on $F_d$; the cusp meridian
has value $p'=\pm1$ under $\Theta$.

Write $N$ for the cusp tube and $M=\partial N$. The exterior-square
matrix of its monodromy has saturated image generated by
$\widehat w\wedge\widehat\delta$ and
$(d\widehat\gamma)\wedge\widehat w+d\widehat u\wedge\widehat\delta$.
Thus the Wang sequence makes $H_2(M)$ free of rank six. Sweeping the
two angular circles $\widehat w,\widehat\delta$ along the toric
meridian gives two classes $\xi_w,\xi_\delta$ which map under the
Wang projection to a basis of the invariant lattice. Both bound
integrally in $N$: sweep the corresponding circle along the extended
zero-section disc in the toric chart. Their span is primitive in
$H_2(M)$. On the other hand
$H_3(N,M)\simeq H^3(N)\simeq H^3(\widetilde W_p)=\mathbb Z^2$;
exactness bounds the kernel of $H_2(M)\to H_2(N)$ by rank two.
Consequently that kernel is exactly
$\mathbb Z\xi_w\oplus\mathbb Z\xi_\delta$.

Cutting $\xi_w$ into the two finite pieces gives connecting boundary
$\pm(-\widehat u-\widehat w+\widehat\delta)$, a primitive generator,
whereas $\xi_\delta$ has zero connecting boundary. The latter belongs
to $\operatorname{coker}\alpha_2$. If $\mathcal A$ is the class dual
to $\Phi$, then $\mathcal A|_{F_d}=2q_d$. Integration along the
free $\delta$-circle on $\widetilde X^\circ$ gives
\[
 (\iota_\delta\mathcal A)|_{F_d}=\pm12d\gamma',
 \qquad \iota_\delta\mathcal A=\pm d\Theta.
\]
The second equality follows from injectivity of restriction in degree
one. Evaluation on the meridian therefore gives
$\Phi(j_*\xi_\delta)=\pm d p'=\pm d$. Hence
\[
 H_2(\widetilde X^\circ)/j_*\ker(H_2(M)\to H_2(N))
 \simeq\mathbb Z/d.
\]
This evaluates the actual swept torus, rather than inferring its
coefficient from rational homology.

Finally, the exact sequence of the pair
$(\widetilde X_p,\widetilde X^\circ)$ and excision give
\[
 0\longrightarrow\mathbb Z/d\longrightarrow H_2(\widetilde X_p)
 \longrightarrow\ker\{H_2(N,M)\to H_1(\widetilde X^\circ)\}
 \longrightarrow0.
\]
Poincare--Lefschetz duality identifies $H_2(N,M)$ with
$H^4(\widetilde W_p)=\mathbb Z^d$. The extended section disc has
boundary the meridian, whose value is $p'=\pm1$, so the final map is
onto $\mathbb Z$ and its kernel is free of rank $d-1$. The sequence
splits as an abelian-group sequence. The rational calculation of
D.1.3 gives $b_3=0$; duality and universal coefficients now yield
$H_3=\operatorname{Ext}(H_2,\mathbb Z)=\mathbb Z/d$,
$H_4=H^2=\mathbb Z^{d-1}$ and $H_5=0$.
\end{proof}
